\documentclass[11pt]{article}
\usepackage[margin=1in]{geometry}
\usepackage{amsmath,amssymb,mathtools}
\usepackage{microtype}
\setlength{\parindent}{0pt}
\setlength{\parskip}{0.7em}

\newcommand{\I}{\mathbb{I}}
\newcommand{\bsigma}{\boldsymbol{\sigma}}

\begin{document}

\begin{center}
  {\Large\bfseries Problem 1 -- Question 2}\\[0.3em]
  {\large Axis-angle representation of an arbitrary \(U(2)\) matrix}
\end{center}

Let \(U\in U(2)\). Since \(\det U\) has unit modulus, choose
\[
 \alpha=\frac{1}{2}\arg(\det U)
\]
and define
\[
 V=e^{-i\alpha}U.
\]
With a consistent choice of the square root of \(\det U\), one has
\(\det V=1\), so \(V\in SU(2)\). Every \(SU(2)\) matrix has the form
\[
 V=
 \begin{pmatrix}
 a&b\\
 -b^*&a^*
 \end{pmatrix},
 \qquad |a|^2+|b|^2=1.
\]
Write
\[
 a=a_R+ia_I,\qquad b=b_R+ib_I.
\]
Then
\[
 V
 =a_R\I+i\left(b_I X+b_RY+a_I Z\right).
\]
The coefficient vector is real, and the normalization condition becomes
\[
 a_R^2+a_I^2+b_R^2+b_I^2=1.
\]

Choose \(\theta\) so that
\[
 \cos\left(\frac{\theta}{2}\right)=a_R,
 \qquad
 s=\sin\left(\frac{\theta}{2}\right)
 =\sqrt{a_I^2+b_R^2+b_I^2}.
\]
If \(s\neq0\), define the real unit vector
\[
 \mathbf n
 =-\frac{1}{s}\left(b_I,b_R,a_I\right).
\]
Using the result of Question 1,
\[
 V
 =\cos\left(\frac{\theta}{2}\right)\I
  -i\sin\left(\frac{\theta}{2}\right)\mathbf n\cdot\bsigma
 =R_{\mathbf n}(\theta).
\]
Therefore
\[
 \boxed{U=e^{i\alpha}R_{\mathbf n}(\theta).}
\]

When \(s=0\), \(V=\I\) or \(V=-\I\). These cases may be represented by
\(\theta=0\) or \(\theta=2\pi\), respectively, with an arbitrary unit axis
\(\mathbf n\). The parameters are not unique: changing the branch of
\(\alpha\), replacing \((\mathbf n,\theta)\) by
\((-\mathbf n,-\theta)\), and using the \(4\pi\)-periodicity of spin
rotations give equivalent representations.

\paragraph{Coordinate-free parameter extraction.}
For \(s\neq0\), the same result can be written as
\[
 \cos\left(\frac{\theta}{2}\right)=\frac{1}{2}\operatorname{Tr}V,
 \qquad
 n_j=\frac{i}{2s}\operatorname{Tr}(\sigma_jV),
 \quad j\in\{X,Y,Z\}.
\]
For \(V\in SU(2)\), these quantities are real.

\end{document}
